\section{Operational Configuration Notes}
\label{app:config}

\sys is environment-configured, and the surface is large --- which
\secref{sec:compare} lists as a genuine onboarding cost. This appendix records the
settings whose \emph{semantics} a reader of this paper needs, rather than a full
inventory.

\para{The two database URLs are independent, deliberately.} \sys's database URL is
configured separately from \mlflow{}'s backend store and is not required to equal
it. The bundled stack uses separate logical databases so the two Alembic histories
never compete for one schema-version row (\secref{sec:impl}). Pointing both at one
database is possible and is a mistake that surfaces later as a migration conflict
rather than immediately as an error.

\para{``Artifact store'' names two different things.} \mlflow{}'s artifact root
supports \mlflow{}'s full backend set, including GCS. \sys's own storage service
accepts \code{local} and \code{s3} only, with MinIO reached through S3
compatibility; anything else is rejected at construction. Reading support on one as
support on the other is a configuration error waiting to happen
(\figref{fig:topologies}).

\para{Background work is gated at two levels.} A master switch enables the loop
lifecycle; three loops additionally have independent enable flags, and one does not.
Disabling the master switch leaves a deployment that serves requests and accumulates
queued work indefinitely --- which is the correct behaviour for a read-only replica
and a silent failure for anything else.

\para{Rate limits are per process.} The failed-login throttle and the API rate limiter
are in-memory token buckets (\secref{sec:exec}). With \mlflow{}'s default of four
gunicorn workers, a configured limit of $L$ yields an effective ceiling near $4L$.
Size the configured value as $L/W$ for $W$ workers, or pin a single worker.

\para{The service body cap is per request.} Published-service invokes count raw
request chunks against a configured maximum, defaulting to 1~MiB. Callers needing to
pass large content should place it in managed storage and pass a reference. Raising
the cap raises a per-request bound and provides no flood protection
(\appref{app:api}).

\para{Provider selection defaults to fakes.} The LLM, evaluation, and promoter
provider selections each default to a deterministic fake, so a fresh deployment
boots and is fully navigable with no API key. The corollary matters for anyone
interpreting a running instance: a deployment that appears to be evaluating
candidates may be running the fake evaluator, and the provider selection should be
checked before any score is believed.

\para{The credential bootstrap is loopback-only.} A convenience default credential is
created only while the account table is empty, only under the loopback launcher, and
never resets an existing password. Any network-reachable deployment must leave the
insecure-default flag disabled and supply a strong password source. The bundled
container stack is a loopback-bound development environment and not a production
template (\secref{sec:impl}).

\para{Single-environment constants are seams, not switches.} The constants that mark
the single-environment boundary are visible in both the frontend and the backend
(D9, \tabref{tab:decisions}). Changing them alone is not a supported activation: a
real multi-environment mode needs configuration, migrations, authorization, and
compatibility tests (\secref{sec:future}).
